\documentclass[10pt]{article}
\usepackage{xeCJK}
\xeCJKsetup{CJKecglue={\hspace{0.1em}}}
\usepackage[a4paper, top=20.4mm, bottom=20.4mm, left=22mm, right=22mm]{geometry}
\usepackage{titlesec}
\usepackage{zhnumber}
\usepackage{graphicx}
\usepackage{setspace} 
\usepackage{color}
\usepackage{subcaption}
\usepackage{tabularx}
\usepackage{changepage}
\usepackage{tabularx, array}
\usepackage{amsmath}
\usepackage{siunitx}
\usepackage{multirow}
\usepackage{diagbox}
\usepackage{enumitem}
\usepackage{extarrows}
\usepackage{bm}
\setlength{\jot}{6pt}
\setlist[enumerate,1]{label=(\arabic)}
\newCJKfontfamily\CJKrm{標楷體}
\newcolumntype{Y}{>{\centering\arraybackslash}X}

\usepackage[colorlinks=true, linkcolor=black, urlcolor=black]{hyperref}

\titleformat{\section}{\Large\bfseries}{\thesection}{0.2em}{}
\titleformat{\subsection}{\large\normalfont}{\thesubsection}{0.5em}{}
\renewcommand{\thesubsection}{\arabic{subsection}.}
\renewcommand{\thesection}{\zhnum{section}、}
\xeCJKsetup{AutoFakeBold=true, AutoFakeSlant=true}
\setCJKmainfont{標楷體}
\begin{document}

原理:
\\
對於一個在理想流體中平移的理想球體，流體力學（勢流理論）嚴格證明其附加質量為排開流體質量的一半：
$$m_a = C_m \rho_w V$$
其中 $C_m$ 為附加質量係數（對於球體 $C_m = 0.5$），$\rho_w$ 為流體（水）密度，$V$ 為體積。

以下推導假設擺錘為完美球體，並處於小角度擺動狀態。

單擺在水中受到的力包含：
 有效重力 (包含浮力修正)： $W_{eff} = (\rho - \rho_w)Vg$
 黏滯阻尼力： 假設為線性阻尼 $F_d = -b v = -b L \dot{\theta}$
 附加質量慣性力： $F_a = -m_a a = -C_m \rho_w V L \ddot{\theta}$

根據牛頓第二運動定律的轉動形式 $\sum \tau = I_0 \ddot{\theta}$，其中本體慣性矩 $I_0 = m L^2 = \rho V L^2$（假設擺線質量不計且視為質點擺）：
$$- (\rho - \rho_w) V g L \sin\theta - b L^2 \dot{\theta} - C_m \rho_w V L^2 \ddot{\theta} = \rho V L^2 \ddot{\theta}$$


將與加速度 $\ddot{\theta}$ 相關的項移至同一側，合併「真實質量」與「附加質量」：
$$(\rho V + C_m \rho_w V) L^2 \ddot{\theta} + b L^2 \dot{\theta} + (\rho - \rho_w) V g L \sin\theta = 0$$

兩邊同除以 $V L^2$ 並使用小角度近似 $\sin\theta \approx \theta$：
$$(\rho + C_m \rho_w) \ddot{\theta} + \frac{b}{V} \dot{\theta} + \frac{(\rho - \rho_w)g}{L} \theta = 0$$

將其化為標準的阻尼震盪方程形式 $\ddot{\theta} + 2\gamma \dot{\theta} + \omega_0'^2 \theta = 0$：
$$\ddot{\theta} + \frac{b}{V(\rho + C_m \rho_w)} \dot{\theta} + \frac{(\rho - \rho_w)g}{(\rho + C_m \rho_w)L} \theta = 0$$

從上式可知，系統在水中的無阻尼自然角頻率 $\omega_0'$ 為：
$$\omega_0' = \sqrt{ \frac{(\rho - \rho_w)g}{(\rho + C_m \rho_w)L} }$$

若測得阻尼常數 $\gamma = \frac{b}{2V(\rho + C_m \rho_w)}$，則實際觀測到的受阻尼擺動角頻率 $\omega_d$ 為：
$$\omega_d = \sqrt{\omega_0'^2 - \gamma^2}$$

在低黏滯流體（如水）且阻尼不大的假設下 ($\gamma \ll \omega_0'$)，水中的擺動週期 $T_{water}$ 近似為：
$$T_{water} \approx \frac{2\pi}{\omega_0'} = 2\pi \sqrt{ \frac{(\rho + C_m \rho_w)L}{(\rho - \rho_w)g} }$$

---



2. 理想無阻尼近似下的 $C_m$ 推導 (Ideal Undamped Approximation)

首先考慮最基礎的情況（忽略流體黏滯性造成的週期偏移）。
在空氣中（假設空氣密度 $\rho_{air} \approx 0$，附加質量效應忽略不計），單擺的自然角頻率 $\omega_{air}$ 與週期 $T_{air}$ 的關係為：
$$\omega_{air}^2 = \frac{g}{L} = \left( \frac{2\pi}{T_{air}} \right)^2$$

在水中，根據我們前述的動力學方程式，系統的無阻尼自然角頻率 $\omega_0'$ 為：
$$\omega_0'^2 = \frac{(\rho - \rho_w)g}{(\rho + C_m \rho_w)L}$$

將兩式相除，消去幾何常數 $g$ 與 $L$，我們可以得到空氣與水中頻率平方的無量綱比值：
$$\frac{\omega_{air}^2}{\omega_0'^2} = \frac{\rho + C_m \rho_w}{\rho - \rho_w}$$

假設流體阻尼極小，可以直接用測量到的週期 $T_{water}$ 近似 $2\pi/\omega_0'$，則上式可寫為週期的平方比：
$$\left( \frac{T_{water}}{T_{air}} \right)^2 = \frac{\rho + C_m \rho_w}{\rho - \rho_w}$$

經過代數移項，將待測目標 $C_m$ 獨立於等式左側，即可得到基礎測量公式：
$$C_m = \left( \frac{\rho}{\rho_w} \right) \left[ \left( \frac{T_{water}}{T_{air}} \right)^2 - 1 \right] - \left( \frac{T_{water}}{T_{air}} \right)^2$$

> 物理意義詮釋：此公式表明，只要我們知道擺錘的密度 $\rho$ 與水的密度 $\rho_w$，並透過感測器精準量測同一個擺錘在空氣中與水中的週期比 $(T_{water}/T_{air})$，就能直接計算出 $C_m$。


4. 驗證與比對：將取得的參數代入精確測量方程式計算 $C_m$，並與勢流理論預測的理論值 $0.5$ 進行誤差分析。我們可進一步改變擺錘材質（改變 $\rho$），繪製 $\frac{\omega_{air}^2}{\omega_d^2 + \gamma^2}$ 對 $\frac{\rho}{\rho-\rho_w}$ 的關係圖，透過線性迴歸的斜率與截距來驗證模型的普適性。
\end{document}